\section{Paired Replay Protocol}
\subsection{Demonstrations, interventions and outcomes}
For each layout, we record a scripted demonstration, restore the initial state and replay to a fixed horizon. Policies receive the source tape and measured hand positions. Gripper commands pass through unchanged. We use object state, fault labels and task predicates to construct and score the experiments. We verify initial states and source tapes across paired policies and retain failed demonstrations.

Fetch push and pick-and-place use Gymnasium-Robotics 1.4.2 and MuJoCo 3.3.7, with a 50-step horizon and 50\,mm per-axis translation scale~\cite{plappert2018fetch}. Meta-World 3.1.1 and MuJoCo 3.3.0 use 10\,mm increments relative to Sawyer's previous mocap target~\cite{yu2020metaworld,mclean2025metaworld}. We use its released push, pick-place and drawer-open experts and evaluate at step 150. Each layout undergoes four conditions: clean replay, five consecutive translations overwritten by zero, a 20\,N hand force along negative world $y$ for three steps, and both interventions. We draw onset times from steps 7--29 on Fetch and 8--65 on Sawyer. All policies in a pair share those times. Fault type and onset are hidden during ordinary replay.

The initial comparison uses 24 held-out layouts per task, or 120 layouts and 3,360 branch executions. Seven clean Fetch push demonstrations miss the task goal and remain in the sample. Follow-up comparisons use separately identified cohorts. We score the native success predicate at the \emph{last} step. For each episode the hand-path RMS error is $\sqrt{T^{-1}\sum_{t=1}^{T}\|x_t-x_t^\star\|_2^2}$. We average this error across episodes. Paired bootstrap intervals resample layouts within each task, carrying all conditions for each layout together. Equal task and intervention weights define this stress test. They have no assumed deployment frequency.

\subsection{Information supplied to each policy}
A command-history policy starts with the known origin and updates it using the commands it sends. An undetected overwrite breaks that recursion. A target read resets the estimate. In the original event policy, two consecutive hand deviations above 5\,mm trigger a read, with a ten-step cooldown and at most two reads per trial. A scheduled policy reads at steps 16 and 32 on Fetch or 50 and 100 on Sawyer. The Sawyer times and $k=0.5$ came from a development comparison of three schedules and two gains. A lower-cost variant checks the residual at the two scheduled times only.

The hand-based observer uses the recorded offset between source hand and source origin. Let $\tilde b_t$ be the live origin predicted from previous estimates and sent commands, and let $b_t^\star$ be the source origin. The correction is
\begin{equation}
 \hat b_t=C\!\left(\tilde b_t+0.25\big[x_t+(b_t^\star-x_t^\star)-\tilde b_t\big]\right).\label{eq:observer}
\end{equation}
It places the recorded offset at the current hand, blends this candidate with the prediction, and propagates the updated estimate using its next submitted command. The read and zero-read branches share this observer, including its gain. We froze it after an eight-layout development trial. It assumes access to a clean source tape, an initialized target and time-aligned hand measurements.

For the limited-resolution comparisons we delay both hand measurements and requested targets by two steps. We propagate a delayed target with submitted commands and transport the last measured hand residual along the demonstration. Only the replay's live hand is rounded to a 1 or 5\,mm grid in world coordinates. The event policy may read once after two steps above 5\,mm. Before testing new layouts at 5\,mm, we fix a read at step 41 (the median development event time) and define a second schedule by permuting the 96 development event times independently of the new motions.
